\section{Baire spaces}\label{sec:baire_spaces}

\begin{definition}\label{def:meager_set}\mcite[def. 2.1]{Rudin1991FunctionalAnalysis}
  Any countable union of \hyperref[def:topologically_dense_set]{nowhere dense sets} is called \term{meager} or a \term{first category set} (see \cref{rem:baire_category_terminology} for terminology). If a set is not meager, we call it \term{nonmeager} or a \term{second category set}.
\end{definition}

\begin{remark}\label{rem:baire_category_terminology}
  \Textcite{Baire1899FonctionsDeVariablesRéelles}, when proving \fullref{thm:baire_category_theorem}, distinguished between \hyperref[def:meager_set]{meager} and nonmeager sets by saying that they are of the \enquote{first category} (\enquote{\textfr{première catégorie}}) and of the \enquote{second category} (\enquote{\textfr{deuxième catégorie}}), correspondingly.

  As described in \cref{rem:category_as_graph}, half a century later, Eilenberg and Mac Lane introduced the unrelated category theory. In the meantime, general topology developed, and Baire's ideas were generalized. Category theory also found their use in topology and even analysis. We will give some examples of how modern authors handle this ambiguity.

  \begin{itemize}
    \item Many authors do not discuss category theory and prefer Baire's terminology, for example
    \cite[159]{Kunen2013SetTheory},
    \cite[def. 25.1]{Willard2004GeneralTopology},
    \cite[132]{Carothers2000RealAnalysis},
    \cite[288]{HasselblattKatok1995DynamicalSystems},
    \cite[def. 3.2.1]{Sally2013Analysis},
    \cite[thm. A.58]{Lee2012SmoothManifolds},
    \cite{Rudin1987RealAndComplexAnalysis} and
    \cite[21]{КанторовичАкилов1984ФункциональныйАнализ}.

    \Textcite[258]{Birkhoff1967LatticeTheory} and \textcite[174]{Folland1999RealAnalysis} also use Baire's terminology in the sections on Baire's theorem, however outside that section they use category theory. Folland states
    \begin{displayquote}
      A more modern and more descriptive synonym of \enquote{of the first category} is \textbf{meager}. The complement of a meager set is called \textbf{residual}.
    \end{displayquote}

    \Textcite[201]{Kelley1975GeneralTopology} does not discuss category theory, but prefers the term \enquote{meager}.

    \item On the other end, some authors use category theory and state Baire's theorem without auxiliary terminology, for example
    \cite[78]{КолмогоровФомин2004ФункциональныйАнализ},
    \cite[30]{ФедорчукФилиппов2006ОбщаяТопология} and
    \cite[162]{Хелемский2014ФункциональныйАнализ}.

    Some authors avoid auxiliary terminology even without using categories, for example
    \cite[thm. 3.9.3]{Engelking1989GeneralTopology} and
    \cite[50]{Deimling1985NonlinearFA}.
  \end{itemize}
\end{remark}

\begin{proposition}\label{thm:def:meager_set}\mcite[43]{Rudin1991FunctionalAnalysis}
  \hyperref[def:meager_set]{Meager sets} have the following basic properties:
  \begin{thmenum}
    \thmitem{thm:def:meager_set/union} A countable union of meager sets is meager.
    \thmitem{thm:def:meager_set/subset} A subset of a meager set is meager.
    \thmitem{thm:def:meager_set/homeomorphism} The \hyperref[def:homeomorphism]{homeomorphic} image of a set \( A \) is meager if and only if \( A \) itself is meager.
  \end{thmenum}
\end{proposition}
\begin{proof}
  \SubProofOf{thm:def:meager_set/union} Follows from \cref{thm:countably_infinite_union_of_countably_infinite_sets}.
  \SubProofOf{thm:def:meager_set/subset} Fix a meager set \( A \) and let \( B \subseteq A \). Then \( A = \bigcup_{k=1}^\infty A_k \) for some nowhere dense sets \( A_1, A_2, \ldots \). By \cref{thm:def:topologically_dense_set/nowhere_dense_subset}, the sets \( A_1 \cap B, A_2 \cap B, \ldots \) are also nowhere dense. But
  \begin{equation*}
    B
    =
    A \cap B
    =
    \left(\bigcup_{k=1}^\infty A_k \right) \cap B
    \reloset {\ref{thm:boolean_algebra_of_subsets}} =
    \bigcup_{k=1}^\infty (A_k \cap B).
  \end{equation*}

  Therefore, \( B \) is also nowhere dense.

  \SubProofOf{thm:def:meager_set/homeomorphism}
\end{proof}

\begin{definition}\label{def:baire_space}
  A topological space is called a \term{Baire space} if any of the following equivalent conditions hold:
  \begin{thmenum}
    \thmitem{def:baire_space/meager} Every nonempty open set is nonmeager.
    \thmitem{def:baire_space/dense} A countable intersection of dense sets is dense.
  \end{thmenum}
\end{definition}
\begin{proof}
  \EquivalenceSubProof{def:baire_space/meager}{def:baire_space/dense} Follows from \cref{thm:def:topologically_dense_set/complement} and \fullref{thm:de_morgans_laws}.
\end{proof}

\begin{proposition}\label{thm:open_subspace_of_baire_space_is_baire}
  Every open subspace of a \hyperref[def:baire_space]{Baire space} is a Baire space.
\end{proposition}
\begin{proof}
  Let \( (X, \mscrT) \) be a Baire space and let \( (X', \mscrT_{X'}) \) be an open \hyperref[def:topological_subspace]{subspace} with the canonical embedding \( \iota: X' \to X \). The proposition holds vacuously if \( X' = \varnothing \), so we assume that \( X' \neq \varnothing \).

  Note that \( \iota \) is continuous by definition, however it is also an open map because if \( U \in T_{X'} \), then \( \iota(U) = U \cap X' \) is open in \( X \) as the intersection of two open sets. Therefore, it is a homeomorphic embedding and by \cref{thm:def:meager_set/homeomorphism}, \( U \) is meager if and only if \( \iota(U) \) is meager. Since \( X \) is a Baire space, \( \iota(U) \) is not meager and hence \( U \) is also not meager.

  We showed that every nonempty open set \( U \in T_{X'} \) is nonmeager, therefore \( X' \) is a Baire space.
\end{proof}

\begin{theorem}[Baire category theorem]\label{thm:baire_category_theorem}\mcite{Rudin1991FunctionalAnalysis}
  \begin{thmenum}
    \thmitem{thm:baire_category_theorem/metric} \hyperref[def:complete_metric_space]{Complete metric spaces} are \hyperref[def:baire_space]{Baire spaces}.
    \thmitem{thm:baire_category_theorem/compact} \hyperref[def:locally_compact_space]{Locally compact} \hyperref[def:separation_axioms/T2]{Hausdorff} spaces are Baire spaces.
  \end{thmenum}
\end{theorem}
